\section{Evaluation}
\label{sec:evaluation}

We evaluate ALS along two dimensions: functional correctness and security
property verification. Five functional (P-KF) and eight non-functional
(P-KNF) test scenarios are executed in the prototype environment described
in Section~\ref{sec:implementation}.

\subsection{Functional Correctness}

\textbf{P-KF-01 (Centralized ingestion).} ALS receives and records events
from all four sources concurrently via \texttt{POST /api/events}. Records
from interleaved sessions are ordered correctly by UUID v7 timestamp;
no record is lost or duplicated.

\textbf{P-KF-02 (Event coverage).} All 13 event types (EV1--EV13) are
triggered by their corresponding DecMed operations. Each appears in the log
file with the correct \texttt{action} label and fully populated
\texttt{AuditEventDetails} fields.

\textbf{P-KF-03 (Schema consistency).} Audit records produced across all
event types parse correctly using the same \texttt{AuditRecord}
deserializer; no field is absent or mistyped.

\textbf{P-KF-04 (Append-only writes).} Log files are opened in append mode;
no in-place modification of existing records is possible through the ALS
API surface.

\textbf{P-KF-05 (Rotation and on-chain commitment).} After rotation, the
IPFS CID and file hash are verifiably committed to the IOTA smart contract
as a frozen object retrievable by any IOTA node.

All five scenarios pass.

\subsection{Security Property Verification}

\textbf{P-KNF-01 (Tamper detection).} Modifying a single byte in any
record in a sealed log file causes the hash-chain verifier to report a
broken chain at that position, with the correct record index identified.

\textbf{P-KNF-02 (Signature enforcement).} Events submitted with an invalid
or absent Ed25519 signature are rejected by the Audit Receiver before
entering the queue; no \texttt{AuditRecord} is created for rejected events.

\textbf{P-KNF-03 (Record standardization).} All records across all event
types conform to the JSON-Lines schema; automated schema validation
reports zero violations.

\textbf{P-KNF-04 (IPFS persistence).} Pinned log files remain retrievable
after triggering a manual garbage-collection run on the IPFS Cluster node;
unpinned files of the same size are removed.

\textbf{P-KNF-05 (On-chain immutability).} Attempts to mutate a frozen
\texttt{RotationRecord} object via the IOTA RPC return a
\texttt{MoveAbort} error; the object is unmodifiable by any transaction.

All eight non-functional scenarios pass.

\subsection{Gap Resolution Traceability}

Table~\ref{tab:traceability} maps each gap identified in
Section~\ref{sec:threat-analysis} to the ALS mechanism that resolves it and
the test that verifies the resolution.

\begin{table}[htbp]
    \caption{Gap--Solution--Test Traceability}
    \label{tab:traceability}
    \begin{center}
        \begin{tabular}{c p{3cm} p{2.5cm} c}
            \toprule
            \textbf{Gap} & \textbf{ALS Mechanism} & \textbf{Test} \\
            \midrule
            M-TA-01 & Single \texttt{POST /api/events} endpoint & P-KF-01 \\
            M-TA-02 & 13-type event catalog from STRIDE & P-KF-02 \\
            M-TA-03 & Hash-chaining + frozen on-chain object & P-KNF-01, P-KNF-05 \\
            M-TA-04 & Uniform \texttt{AuditRecord} JSON-Lines schema & P-KNF-03 \\
            M-TA-05 & IPFS pinning + IOTA frozen object & P-KNF-04, P-KNF-05 \\
            \bottomrule
        \end{tabular}
    \end{center}
\end{table}

\subsection{Limitations}

The evaluation was conducted on IOTA Localnet rather than Testnet or Mainnet;
behavior under production network conditions---particularly transaction
latency variability and IPFS Cluster node failures---has not been measured.
The rotation interval was set to 60 seconds for testing; production
deployments should use 300 seconds or longer to reduce on-chain transaction
overhead. Ed25519 key pairs are currently ephemeral, regenerated at each
component restart, which limits forensic continuity across process restarts;
a persistent key management solution (e.g., HSM or secure enclave) is
required for full production-grade non-repudiation.